\begin{abstract}
\noindent\textbf{Introduction:} Recent analyses of the US Food and Drug Administration Adverse Event
Reporting System (FAERS) reported many drug--hearing-loss signals absent from labels. We assessed which survive adjustment for
hearing-specific biases and external replication.

\noindent\textbf{Methods:} Disproportionality analysis of 19.9 million deduplicated FAERS reports
(2004--2026) and 1.2 million Canada Vigilance reports, following a protocol committed before
analysis. For hearing loss and tinnitus, crude reporting odds ratios (signal: lower 95\% bound $>1$, $\geq$3 cases) were
compared with Mantel--Haenszel estimates stratified by age, sex and period after excluding ear-related indications, reports
naming established ototoxic drugs and lawyer reports. Methods were compared on a label-derived reference set (area under the
receiver operating characteristic curve, AUC) and by replication in Canada Vigilance; robust unlabelled signals underwent
case-by-case review.

\noindent\textbf{Results:} Adjustment changed the hearing-loss signal list (63 of
300 crude signals lost) but not agreement with labels (AUC, crude vs adjusted lower bound:
0.718 vs 0.723). Unlabelled signals replicated in Canada more than background but less than labelled ones; power-calibrated (post hoc), labelled signals replicated as predicted (observed/expected
1.02, 95\% CI 0.77--1.32)
and unlabelled ones did not (0.44, 0.29--0.62),
also within strata of consumer reporting. Solicited reports, report clusters and shared
co-medication produced spurious replications. Of 22 robust unlabelled signals (post hoc), 15 had artefactual or known
explanations; 7 remained (2--9 under
alternative choices), led by pregabalin.

\noindent\textbf{Conclusions:} Most unlabelled FAERS hearing-loss signals do not survive adjusted replication and provenance
review. Findings are hypothesis-generating; cross-database replication needs adjusted estimates and explicit
handling of solicited reports.
\end{abstract}